\documentclass[11pt]{article}

% Shared preamble for the SKJS soundness corpus (main paper + deep dives).
% Loaded via \input in each document's preamble. Documents may define
% additional packages, theorem environments, and macros AFTER this input,
% but must not redefine anything here.
\usepackage{amsmath}
\usepackage{amssymb}
\usepackage{amsthm}
\usepackage{mathtools}
\usepackage{stmaryrd}
\usepackage[margin=1in]{geometry}
\usepackage{tcolorbox}
\usepackage{booktabs}
\usepackage{array}
\usepackage{longtable}
\usepackage[hyphens]{url}
\urlstyle{tt}

% Long typewriter tokens (paths, identifiers) defeat hyphenation;
% allow inter-word stretch instead of overfull lines.
\emergencystretch=2em

\newtheorem{theorem}{Theorem}
\newtheorem{lemma}[theorem]{Lemma}
\newtheorem{definition}[theorem]{Definition}
\newtheorem{corollary}[theorem]{Corollary}
\newtheorem{axiom}[theorem]{Axiom}
\newtheorem{obligation}[theorem]{Obligation}

% Summary box for key definitions
\newtcolorbox{keydef}{
  colback=gray!5,
  colframe=gray!50,
  boxrule=0.5pt,
  arc=2pt,
  left=4pt,right=4pt,top=2pt,bottom=2pt
}

% Core semantic notation shared by every document
\newcommand{\sem}[1]{\llbracket #1 \rrbracket}
\newcommand{\Val}{\mathsf{Val}}

% Implementation anchor: file / class / function at the pinned revision
\newcommand{\impl}[1]{\par\noindent{\footnotesize\textsf{impl:} \texttt{#1}}}

% Single source for the modeled revision. Every document's \date uses this;
% the ledger and tooling read the same values.
\newcommand{\skjspinhash}{9108f339}
\newcommand{\skjspindate}{2026-07-02}
\newcommand{\skjspin}{Models revision \texttt{\skjspinhash} of \texttt{main} (\skjspindate)}



\newtheorem{finding}[theorem]{Finding}

% Long typewriter tokens (spec names, file paths) defeat hyphenation;
% allow a little inter-word stretch instead of overfull lines.
\emergencystretch=2em

% Macros copied from the main paper (skjs-type-soundness.tex, rev 9108f339)
\newcommand{\ObjVal}{\mathsf{ObjVal}}
\newcommand{\String}{\mathsf{String}}
\newcommand{\SymVal}{\mathsf{SymVal}}
\newcommand{\Tags}{\mathsf{Tags}}
\newcommand{\Payload}{\mathsf{Payload}}
\newcommand{\rtag}{\mathsf{rtag}}
\newcommand{\tagfn}{\mathsf{tag}}

% Local to this deep dive
\newcommand{\ESVal}{\mathsf{ESVal}}
\newcommand{\estag}{\mathsf{estag}}
\newcommand{\realize}{\gamma}
\newcommand{\ecma}[1]{ECMA-262 \S#1}
\newcommand{\intm}[1]{[[#1]]}

\title{Runtime Grounding:\\
ECMAScript Semantics for the SKJS Value Domain\\
{\large Deep Dive A for ``Refinement and Inference Growth in SKJS''}}
\author{}
\date{\skjspin}

\begin{document}

\maketitle

\begin{abstract}
The main soundness paper models the runtime as an abstract value domain
$\Val$ equipped with predicates ($\approx$, $\mathsf{truthy}$,
$\mathsf{callable}$, $\mathsf{constructible}$, $\mathsf{isArrayV}$,
$P_{\mathsf{inst}}$); three of its ledger rows rest entirely on that
domain: \textbf{L-TAG}, \textbf{L-INST}, and \textbf{L-ISARR}. This
paper is the evidence behind their current table statuses.
SKJS is a type checker only — there is no SKJS code generator, and checked
programs run \emph{unmodified} on a JavaScript engine — so each of these
rows is really a claim about ECMAScript. This deep dive supplies the
grounding: a realization map from the model's value domain onto ECMAScript
language values under which the model's predicates are exactly the
pullbacks of the spec's \texttt{IsStrictlyEqual}, \texttt{ToBoolean},
\texttt{typeof}, \texttt{IsCallable}/\texttt{IsConstructor}, and
\texttt{IsArray}. Under it, \textbf{L-TAG} and \textbf{L-ISARR} are
discharged (\textsc{proved}, modulo two named environment assumptions),
and \textbf{L-INST} is decomposed into sharper named assumptions — two of
which turn out to be \textsc{violated}, with checked-and-executed
counterexamples produced while attempting the discharge. A third
by-product refutes an implementation premise behind \textsc{R-Typeof}:
class constructors implement $\intm{Call}$, but the checker types them as
definitely-not-callable. Section~\ref{sec:status} gives the precise
ledger edits for the main paper.
\end{abstract}

\tableofcontents

%==========================================================================
\section{Purpose and Relation to the Main Paper}
\label{sec:purpose}

This document is deep dive A of the series supporting
\emph{Refinement and Inference Growth in SKJS: a Soundness Model of the
Type Checker} (the \emph{main paper}), pinned like it to revision
\texttt{9108f339}. Its job is to discharge, sharpen, or refute the rows of
the main paper's assumption ledger (its Section ``Assumptions and Proof
Obligations'') that concern the \emph{runtime} rather than the checker:

\begin{itemize}
\item \textbf{L-TAG} — every value carries exactly one
  tag, consistently with its static tag. \emph{SKJS emits no code}: a
  checked program is executed, byte for byte, by an ordinary JavaScript
  engine, so the runtime counterpart of the model's Lemmas A0/A1 is a
  statement about ECMAScript's language types, provable from the
  specification (Section~\ref{sec:tags}) — which is why the ledger anchors
  this row to ECMAScript rather than to any SKJS component.
\item \textbf{L-INST} —
  $P_{\mathsf{inst}}(C) \subseteq \sem{T_C}$ for the instance type built at
  an \texttt{instanceof} narrowing site (Section~\ref{sec:inst}).
\item \textbf{L-ISARR} —
  $\mathsf{isArrayV}(o)$ implies
  $o \in \sem{\mathsf{ReadonlyArray}\langle\top\rangle}$
  (Section~\ref{sec:isarr}).
\end{itemize}

In addition it grounds two definitions the main paper introduces without
external justification and then uses everywhere: the strict-equality
relation $\approx$ (its Definition ``Strict Equality''; used by \textsc{R-Eq},
\textsc{R-EqField}, and constant removal) and
the per-tag truthiness table (its Definition ``Truthiness''; used by
\textsc{R-Truthy}/\textsc{R-TruthyField} and the residue lemma). Both are
shown to be \emph{exactly} the ECMAScript algorithms, transported along
the realization map (Sections~\ref{sec:seq} and~\ref{sec:truthy}).

Failed discharge attempts are reported as numbered
\emph{Findings}. Three of them come with programs that the pinned checker
accepts (verified by running \texttt{skjs check} at the pinned revision)
and that misbehave when executed on Node: they are witnesses, not
conjectures. All spec citations are to ECMA-262 as of the 2025 edition;
clause \emph{numbers} drift between editions, so each citation also names
the clause.

\paragraph{Notation.}
We reuse the main paper's notation unchanged: $\Val$, $\Tags$, $\rtag$,
payloads $(\Payload_k)_k$, $\ObjVal$, the object interface
($\mathsf{props}$, $\mathsf{get}$, $\mathsf{has}$, $\mathsf{callable}$,
$\mathsf{constructible}$, $\mathsf{isArrayV}$), $\approx$,
$\mathsf{truthy}$, and $\sem{\cdot}$. Implementation anchors
(\textsf{impl:}~\texttt{\dots}) point into revision \texttt{9108f339}.

%==========================================================================
\section{Minimal Recap of the Modeled Runtime}
\label{sec:recap}

From the main paper's Section 2 (Runtime Model): $\Tags$ is the
nine-element set
\[
\Tags = \{\mathsf{undef}, \mathsf{void}, \mathsf{null}, \mathsf{bool},
\mathsf{num}, \mathsf{str}, \mathsf{bigint}, \mathsf{sym},
\mathsf{obj}\},
\]
and $\Val = \{(k,p) \mid k \in \Tags, p \in \Payload_k\}$ with
$\Payload_{\mathsf{num}} = \mathbb{F}$ (IEEE-754 doubles including $-0$
and $\mathsf{NaN}$), $\Payload_{\mathsf{bigint}} = \mathbb{Z}$, and
$\Payload_{\mathsf{undef}} = \Payload_{\mathsf{void}} = \{\star\}$. The
$\mathsf{void}$ tag is declared there to be ``a checker fiction for the
runtime value \texttt{undefined}''; this paper makes that statement
precise. The relation $\approx$ is a partial equivalence
($\mathsf{NaN}$ irreflexive, $-0 \approx 0$, the
$\mathsf{undef}/\mathsf{void}$ pair identified);
$\mathsf{truthy}$ is given by per-tag falsy sets; objects come with the
abstract interface above, read as a snapshot (\textbf{L-PURE}).

The three refinement rules whose soundness proofs consume the assumed
rows are \textsc{R-Typeof} (via the $P_{\mathsf{typeof}=k}$ sets and the
$\mathsf{constructible} \implies \mathsf{callable}$ interface fact),
\textsc{R-Inst} (via \textbf{L-INST}), and \textsc{R-IsArray} (via
\textbf{L-ISARR}).

%==========================================================================
\section{The Realization Map: Tags vs.\ ECMAScript Language Types}
\label{sec:tags}

ECMAScript defines exactly eight \emph{language types}: Undefined,
Null, Boolean, String, Symbol, Number, BigInt, and Object
(\ecma{6.1}, ``ECMAScript Language Types''). Every language value belongs
to exactly one of them — the types are introduced as a partition of the
value space, not as overlapping classifications. Write $\ESVal$ for the
set of language values and $\estag(v)$ for the language type of
$v \in \ESVal$.

\begin{definition}[Realization Map]
\label{def:gamma}
$\realize : \Val \to \ESVal$ sends a model value to the language value it
stands for:
\begin{itemize}
\item $\realize(\mathsf{undef}, \star) = \realize(\mathsf{void}, \star)
  = \texttt{undefined}$ and $\realize(\mathsf{null}, \star) =
  \texttt{null}$;
\item on $\mathsf{bool}$, $\mathsf{num}$, $\mathsf{str}$,
  $\mathsf{bigint}$: the evident payload realization — the model chose
  its payload sets to \emph{be} the spec's value sets
  ($\mathbb{F}$ is exactly the Number type, ``double-precision 64-bit
  binary format IEEE 754-2019'' including $-0$ and a single
  $\mathsf{NaN}$, \ecma{6.1.6.1}; $\mathbb{Z}$ is exactly the BigInt
  type, \ecma{6.1.6.2}; $\String$ realizes as UTF-16 code unit
  sequences, \ecma{6.1.4});
\item on $\mathsf{sym}$ and $\mathsf{obj}$: a fixed bijection between the
  abstract payload sets $\SymVal$, $\Payload'_{\mathsf{obj}}$ and the
  runtime symbol/object identities of the executing program.
\end{itemize}
\end{definition}

\begin{theorem}[Grounding of Tags; discharges \textbf{L-TAG}]
\label{thm:tags}
$\realize$ is surjective onto the language values reachable by the
program, and:
\begin{itemize}
\item[(G1)] $\estag(\realize(v))$ is determined by $\rtag(v)$, via the
  evident bijection
  $\Tags \setminus \{\mathsf{void}\} \leftrightarrow$
  \{Undefined, Null, Boolean, Number, String, BigInt, Symbol, Object\},
  extended by $\mathsf{void} \mapsto$ Undefined.
\item[(G2)] $\realize$ is injective except on the single fiber
  $\{(\mathsf{undef},\star), (\mathsf{void},\star)\}$, whose two elements
  realize to the same value \texttt{undefined}.
\item[(G3)] Distinct tags other than the pair
  $\{\mathsf{undef}, \mathsf{void}\}$ realize to disjoint sets of
  language values.
\end{itemize}
Consequently the runtime counterparts of the main paper's Lemmas A0 (tag
consistency) and A1 (tag disjointness) hold: a language value inhabits
denotations of at most one tag, except that \texttt{undefined} inhabits
both $\mathsf{undef}$- and $\mathsf{void}$-tagged denotations.
\end{theorem}
\begin{proof}
(G1) and (G2) are immediate from Definition~\ref{def:gamma}. (G3) follows
from (G1)+(G2) because the eight language types partition $\ESVal$
(\ecma{6.1}). Surjectivity: every language type is covered, and the
payload realizations are onto their value sets (bijections or identity).
\end{proof}

The residual obligation created by (G2) is that no rule of the checker
may depend on \emph{runtime-distinguishing} the tags $\mathsf{undef}$ and
$\mathsf{void}$ — they are indistinguishable after realization. The main
paper already maintains this discipline explicitly: \textsc{R-EqUndef}
and the \texttt{typeof "undefined"} case of \textsc{R-Typeof} treat the
two tags jointly; the truthiness residue keeps $\mathsf{void}$ on both
branches; the one corner where a single tag of the pair could be consumed
alone is closed there by the checker's rejection of
\texttt{void}-typed strict-equality operands. Nothing further is needed
here.

\paragraph{Anchoring.}
There is no code generator anywhere in the pinned revision; the runtime
is whatever engine executes the unmodified source, so the ledger's L-TAG
row is anchored to this theorem (equivalently, to \ecma{6.1}), not to any
SKJS component. Note the row's phrase ``consistent with its static tag''
asks, at full strength, that checked expressions inhabit their checked
types; that is \textbf{L-TC}, a different row, and this paper
deliberately does not fold it back in.

\subsection{The \texttt{typeof} table}
\label{sec:typeof-table}

The \texttt{typeof} operator (\ecma{13.5.3}, ``The \texttt{typeof}
Operator'', Table ``typeof Operator Results'') maps every language value
to a string as follows, where \emph{callable} means ``implements the
$\intm{Call}$ internal method'':

\begin{center}
\renewcommand{\arraystretch}{1.15}
\begin{tabular}{@{}lll@{}}
\toprule
Language type of $v$ & \texttt{typeof} $v$ & model tag(s) \\
\midrule
Undefined & \texttt{"undefined"} & $\mathsf{undef}$, $\mathsf{void}$ \\
Null & \texttt{"object"} & $\mathsf{null}$ \\
Boolean & \texttt{"boolean"} & $\mathsf{bool}$ \\
Number & \texttt{"number"} & $\mathsf{num}$ \\
String & \texttt{"string"} & $\mathsf{str}$ \\
BigInt & \texttt{"bigint"} & $\mathsf{bigint}$ \\
Symbol & \texttt{"symbol"} & $\mathsf{sym}$ \\
Object, not callable & \texttt{"object"} & $\mathsf{obj}$ \\
Object, callable & \texttt{"function"} & $\mathsf{obj}$ \\
\bottomrule
\end{tabular}
\end{center}

Two quirks are load-bearing for the checker and both are already encoded
in the main paper's $P_{\mathsf{typeof}=k}$ sets (its Definition
``Evaluation of Condition Forms''):
\texttt{typeof null === "object"} (the \texttt{null} row above; the model
includes $\sem{\texttt{null}}$ in $P_{\mathsf{typeof}=\texttt{"object"}}$),
and the callable/non-callable split of the Object type (the model's
$P_{\mathsf{typeof}=\texttt{"function"}} = \{o \in \ObjVal \mid
\mathsf{callable}(o)\}$). Granting Definition~\ref{def:callable} below
(callable $=$ implements $\intm{Call}$), the pullback along $\realize$ of
``\texttt{typeof} $x$ $=$ $k$'' is \emph{exactly}
$P_{\mathsf{typeof}=k}$ for all eight operands. So the evaluation
semantics of \textsc{R-Typeof} is grounded; the soundness of its
narrowing given that semantics is the main paper's Lemma
``\textsc{R-Typeof} is Branch-Sound'' and is not
re-proved here. (Whether the checker's \emph{types} populate the callable
split correctly is a separate question — see
Finding~\ref{find:ctor-call}.)

\paragraph{Named environment caveat (A-DDA).}
Annex B of the spec (\ecma{B.3.6}, ``The $\intm{IsHTMLDDA}$ Internal
Slot'') sanctions host objects — in practice the single browser object
\texttt{document.all} — for which \texttt{typeof} returns
\texttt{"undefined"} and \texttt{ToBoolean} returns \textsf{false}
despite being objects. Such a value would break rows one and eight of the
table and the truthiness grounding below. We record
\textbf{A-DDA}: \emph{no $\intm{IsHTMLDDA}$ object flows into checked
code}. This holds in any non-browser embedding (Node has none), and the
pinned import stubs describe a Node-style environment
(\texttt{process}, \texttt{fetch}/undici — no DOM).
\impl{import-stubs/@types/skjs/index.d.ts}

\subsection{Where the \texttt{void} tag fits}

The checker distinguishes \texttt{void} from \texttt{undefined} to
restrict \emph{programs} (a \texttt{void}-typed result should not be
consumed), not to describe a runtime distinction; Theorem~\ref{thm:tags}
(G2) is the formal statement that no runtime distinction exists. The
grounded reading of every $\mathsf{void}$-related clause in the main
paper is obtained by applying $\realize$: e.g.\ the cross-tag arm of
$\approx$ (next section) is not a modeling convenience but a
\emph{theorem forced by} $\realize(\mathsf{undef},\star) =
\realize(\mathsf{void},\star)$.

%==========================================================================
\section{Strict Equality: $\approx$ is the Pullback of
\texttt{IsStrictlyEqual}}
\label{sec:seq}

The spec's strict-equality algorithm (\ecma{7.2.16},
\texttt{IsStrictlyEqual}$(x, y)$; the semantics of \texttt{===} via
\ecma{13.11}) is:
(1) if $\estag(x) \neq \estag(y)$, return \textsf{false};
(2) if $x$ is a Number, return
\texttt{Number::equal}$(x,y)$ — which returns \textsf{false} if either
argument is $\mathsf{NaN}$, \textsf{true} if both are $+0$/$-0$ in any
combination, and value equality otherwise (\ecma{6.1.6.1.13});
(3) otherwise return \texttt{SameValueNonNumber}$(x,y)$ — payload
equality per type: trivially \textsf{true} for Undefined and Null,
mathematical-value equality for BigInt, code-unit equality for String,
identity for Boolean/Symbol/Object (\ecma{7.2.12}).

\begin{theorem}[Grounding of $\approx$]
\label{thm:seq}
For all $v, w \in \Val$:
$v \approx w$ iff
$\texttt{IsStrictlyEqual}(\realize(v), \realize(w)) = \textsf{true}$.
\end{theorem}
\begin{proof}
Case on the tag pair. If $\rtag(v), \rtag(w) \in
\{\mathsf{undef},\mathsf{void}\}$: both sides realize to
\texttt{undefined}; step (1) passes, step (3) returns \textsf{true}; and
$v \approx w$ by the first arm of the definition — both sides hold. If
the tags are equal and not in the pair: realization preserves the
language type, step (1) passes, and steps (2)/(3) compute exactly the
payload condition of $\approx$'s second and third arms
($\mathsf{NaN}$-irreflexivity and $-0 = 0$ from
\texttt{Number::equal}; payload/identity equality from
\texttt{SameValueNonNumber} — for BigInt this is mathematical equality
and $\Payload_{\mathsf{bigint}} = \mathbb{Z}$ has no signed zero, so no
analogue of the $-0$ clause arises). If the tags differ and are not the
$\{\mathsf{undef},\mathsf{void}\}$ pair: by Theorem~\ref{thm:tags} (G3)
the realizations have different language types, step (1) returns
\textsf{false}, and $\approx$ does not relate $v, w$ — both sides fail.
\end{proof}

Every property of $\approx$ that the main paper proves or uses —
the partial-equivalence lemma, $\mathsf{NaN} \not\approx \mathsf{NaN}$,
the identification of the $\mathsf{undef}/\mathsf{void}$ pair — is
therefore inherited from the spec rather than stipulated.

\paragraph{Contrast: \texttt{SameValue} and \texttt{Object.is}.}
The spec has two finer equalities: \texttt{SameValue}
(\ecma{7.2.10}; surfaced as \texttt{Object.is}) makes $\mathsf{NaN}$
reflexive and \emph{separates} $+0$ from $-0$; \texttt{SameValueZero}
(\ecma{7.2.11}; used by \texttt{Array.prototype.includes},
\texttt{Map}/\texttt{Set} keying) makes $\mathsf{NaN}$ reflexive and
identifies the zeros. Neither is $\approx$. This matters in two
directions: (i) the checker narrows only on \texttt{===}/\texttt{!==},
so only \texttt{IsStrictlyEqual} needs grounding —
\texttt{Object.is} is not even in the pinned stubs, and
\texttt{Map}/\texttt{Set}/\texttt{includes} are opaque signatures (a subtyping-surface concern); (ii) inside \texttt{OrdinaryHasInstance} the spec
compares prototype-chain entries with \texttt{SameValue}, but on Objects
\texttt{SameValue} and \texttt{IsStrictlyEqual} coincide (both are
identity), so no discrepancy reaches Section~\ref{sec:inst}.

\begin{corollary}[The literal-type closure choice is forced]
\label{cor:zero}
The main paper defines literal-type denotations as same-tag
$\approx$-closures, so the numeric literal type \texttt{0} denotes
$\{+0, -0\}$. Under Theorem~\ref{thm:seq} this is the unique sound
choice for \textsc{R-Eq}'s false branch: \texttt{x !== 0} does not
exclude $x = -0$ at runtime ($-0$ \texttt{=== 0} is \textsf{true}), so a
denotation of \texttt{0} omitting $-0$ would let constant removal
subtract a value still possible. It also matches \texttt{tsc}, which has
no $-0$ literal type: \texttt{const x: 0 = -0} is accepted there, i.e.\
\texttt{tsc}'s literal type \texttt{0} likewise contains both zeros.
\end{corollary}

%==========================================================================
\section{Truthiness: the Falsy Sets are \texttt{ToBoolean}}
\label{sec:truthy}

\texttt{ToBoolean} (\ecma{7.1.2}) is defined by a per-type table:
\texttt{undefined} and \texttt{null} $\mapsto$ \textsf{false}; Boolean
$\mapsto$ itself; Number $\mapsto$ \textsf{false} iff the value is
$+0$, $-0$, or $\mathsf{NaN}$; String $\mapsto$ \textsf{false} iff empty;
BigInt $\mapsto$ \textsf{false} iff $0$; Symbol $\mapsto$ \textsf{true};
Object $\mapsto$ \textsf{true} (with the Annex-B
$\intm{IsHTMLDDA}$ exception covered by \textbf{A-DDA}).

\begin{theorem}[Grounding of Truthiness]
\label{thm:truthy}
Under \textbf{A-DDA}, for all $v \in \Val$:
\[
\mathsf{truthy}(v) = \texttt{ToBoolean}(\realize(v)).
\]
\end{theorem}
\begin{proof}
Row-by-row comparison of the main paper's Definition ``Truthiness'' with
the table
above, through the tag correspondence of Theorem~\ref{thm:tags}: the
$\mathsf{undef}$, $\mathsf{void}$ rows both realize to \texttt{undefined}
(falsy — the model marks both tags falsy, consistently); the numeric
falsy set $\{0, -0, \mathsf{NaN}\}$, the empty string, the zero bigint,
and the always-truthy symbol/object rows are verbatim. The comparison is
exact and complete: neither table has a row the other lacks.
\end{proof}

This grounds the predicate. Note where the burden then sits: the residue
lemma of the main paper (its Lemma ``Residues are Sound'') additionally
needs the
\emph{implementation's lexeme-level} falsiness test to agree with
$\mathsf{truthy}$ on literal spellings — a checker-internal (not
runtime) lexeme-level matter, which remains \textsc{violated} by the
\#655 zero spellings, and is out of scope here.
The runtime side contributes no further caveat.
\impl{TType.sk / TPrimValueKind::isFalsy, SType::falsiness}

%==========================================================================
\section{Callable and Constructible}
\label{sec:call}

\begin{definition}[Grounded predicates]
\label{def:callable}
For $o \in \ObjVal$:
$\mathsf{callable}(o)$ iff $\realize(o)$ implements the $\intm{Call}$
internal method (the spec's \texttt{IsCallable}, \ecma{7.2.3});
$\mathsf{constructible}(o)$ iff $\realize(o)$ implements
$\intm{Construct}$ (\texttt{IsConstructor}, \ecma{7.2.4}).
\end{definition}

``Callable'' therefore means \emph{has the internal method}, not ``can be
invoked without throwing''. This distinction is exactly what the
\texttt{typeof} table keys on, and it is where class constructors live:
a class constructor is an ECMAScript function object whose $\intm{Call}$
\emph{exists but throws} — step ``If $F$.$\intm{IsClassConstructor}$ is
\textsf{true}, throw a \texttt{TypeError}'' in the $\intm{Call}$ of
ECMAScript function objects (\ecma{10.2.1}). Hence
\texttt{typeof (class \{\})} $=$ \texttt{"function"}.

\begin{theorem}[Interface fact:
$\mathsf{constructible} \implies \mathsf{callable}$]
\label{thm:ctor-call}
For every $o \in \ObjVal$, $\mathsf{constructible}(o)$ implies
$\mathsf{callable}(o)$.
\end{theorem}
\begin{proof}
The spec states it directly in the definition of the Object type
(\ecma{6.1.7.2}, ``Object Internal Methods and Internal Slots''):
``every object that supports $\intm{Construct}$ must support
$\intm{Call}$; that is, every constructor must be a function object''
(reiterated among the invariants of the essential internal methods,
\ecma{6.1.7.3}). It is also easy to audit exhaustively across the
function kinds the spec defines:
\emph{ordinary functions} — $\intm{Construct}$ is installed only on
function objects (arrow functions and methods get $\intm{Call}$ only);
\emph{class constructors} — both present, $\intm{Call}$ throwing
(\ecma{10.2.1});
\emph{bound functions} (\ecma{10.4.1}) — always have $\intm{Call}$, and
$\intm{Construct}$ only if the bound target has it;
\emph{Proxy exotic objects} (\ecma{10.5.12}--\S10.5.13) — a proxy has
$\intm{Call}$ iff its target does and $\intm{Construct}$ iff its target
does, so the implication lifts through proxies by induction on the
target chain.
\end{proof}

This discharges the ``JS fact'' quoted in the main paper's Object
Interface (its Definition ``Object Interface'') and used in the
\texttt{"object"} case of its Lemma ``\textsc{R-Typeof} is Branch-Sound'':
status \textsc{proved}.

\begin{finding}[Class constructors refute the checker's
$\mathsf{call} = \mathsf{None}$ typing; unsound \texttt{typeof}
narrowing]
\label{find:ctor-call}
Attempting to ground \textsc{R-Typeof}'s object-branch treatment
end-to-end turned up a genuine unsoundness. The checker's static-side
object type for a class $C$ (the type of the constructor \emph{value})
is built with a construct signature but \emph{no} call signature —
$\mathsf{new} = \mathsf{Some}(\ldots)$,
$\mathsf{call} = \mathsf{None}$
(\texttt{TEnv::newStaticObject} passes \texttt{new} and omits
\texttt{call}). Under the main paper's object semantics,
$\mathsf{call} = \mathsf{None}$ asserts $\neg\mathsf{callable}$ of every
member; but by Definition~\ref{def:callable} and \ecma{10.2.1} the
runtime constructor value \emph{is} callable
(\texttt{typeof C === "function"}). So $\sem{\texttt{typeof } C}$
excludes the value $C$ itself — \textbf{L-TC} is violated at
class-constructor values — and \textsc{R-Typeof}, whose narrowing is
sound \emph{given} membership, misfires. Checked witness (accepted by
the pinned checker, exit 0; \texttt{tsc} rejects it):
\begin{verbatim}
class C { a: number = 1 }
function g(): string {
  const F = C;
  if (typeof F === "function") {
    const s: string = F;   // accepted: the branch narrowed F to bottom
    return s;
  }
  return "not a function";
}
console.log(g());          // Node prints: [class C]
\end{verbatim}
At runtime the true branch is taken; the checker had dropped the sole
object member of $F$'s type (call mode $\mathsf{None}$ $\to$ drop in the
\texttt{"function"}/true arm), leaving an uninhabited type from which
the \texttt{string} annotation is vacuously reachable. $g()$ returns a
class constructor as a \texttt{string}. The false branch of
\texttt{typeof F === "object"} misfires symmetrically. The defect is not
in \textsc{R-Typeof}'s rule (the main paper's proof is correct given
membership) nor in the runtime facts (grounded above); it is in the
static-side type: $\mathsf{call} = \mathsf{None}$ (``definitely not
callable'') should be the maybe-callable mode
$(\mathsf{true}, \sigma)$, or the \texttt{"function"} arm should treat
$\mathsf{new} \neq \mathsf{None}$ as callable, per
Theorem~\ref{thm:ctor-call}. Proposed ledger row:
\textbf{A-CTOR-CALL}, status \textsc{violated}
(Section~\ref{sec:status}); a pinned test is recommended.
\impl{TEnv.sk / newStaticObject (call omitted);
TypeCheckRefine.sk / evalTypeofEquality}
\end{finding}

%==========================================================================
\section{\texttt{instanceof} and \textbf{L-INST}}
\label{sec:inst}

\subsection{The spec algorithms}

Evaluation of \texttt{V instanceof T} calls
\texttt{InstanceofOperator}$(V, T)$ (\ecma{13.10.2}):
\begin{enumerate}
\item if $T$ is not an Object, throw;
\item let $h = \texttt{GetMethod}(T, \texttt{@@hasInstance})$; if $h$
  is not \texttt{undefined}, call $h$ with receiver $T$ and argument
  $V$, and return the \texttt{ToBoolean} of the result;
\item otherwise, if $T$ is not callable, throw; return
  \texttt{OrdinaryHasInstance}$(T, V)$.
\end{enumerate}
The default \texttt{@@hasInstance} method, on
\texttt{Function.prototype}
(\ecma{20.2.3.6}), exists on every function and simply runs
\texttt{OrdinaryHasInstance}; it is non-writable and non-configurable,
but a class can \emph{shadow} it with a static method, and any object can
carry an own \texttt{@@hasInstance}. \texttt{OrdinaryHasInstance}$(C,O)$
(\ecma{7.3.22}): if $C$ is a bound function, recurse on its target; if
$O$ is not an Object, return \textsf{false}; let
$P = \texttt{Get}(C, \texttt{"prototype"})$ (throw if not an Object);
walk $O$'s $\intm{GetPrototypeOf}$ chain, returning \textsf{true} iff
some entry is \texttt{SameValue}-identical to $P$.

Ground $P_{\mathsf{inst}}(C) \triangleq \{v \in \Val \mid
\texttt{InstanceofOperator}(\realize(v), \realize(C)) =
\textsf{true}\}$; the throwing cases satisfy neither branch relation
(main paper, Definition ``Evaluation of Condition Forms''), so they are
vacuous for narrowing. That
$P_{\mathsf{inst}}(C) \subseteq \ObjVal$ is immediate from step 3 of
\texttt{OrdinaryHasInstance} — \emph{provided} evaluation reaches it,
which step 2 does not guarantee. This is the first caveat.

\subsection{What the checker builds, and the gate}

\textsc{R-Inst} fires only on the true branch, and only when the RHS is a
bare name that resolves to a class or interface \emph{definition}
(\texttt{hasClassDef}); a merely Function-typed RHS narrows nothing.
The instance type $T_C$ is built by \texttt{newObject} with type
arguments from \texttt{makeTArgs} — verified at the pin:
\texttt{makeTArgs} generates one \emph{fresh ordinary inference variable}
(\texttt{Subst::gen}, a candidate \texttt{TVID}) per class parameter and
attaches exactly the parameter's \emph{declared constraint}, if any
(\texttt{tvar.setConstraint}); \texttt{Array} is special-cased to the
single argument $\top$.
\impl{TypeCheckRefine.sk / evalCond (InstanceOf arm), evalInstanceof}
\impl{TEnv.sk / hasClassDef, makeTArgs, newObject}

This pins the accurate discharge premise, recorded in the main paper's
L-INST row: the chosen arguments are fresh, but \emph{not} unconstrained
when the class declares parameter bounds
(\texttt{class C<T extends U>}). Freshness plus declared-bounds-only is
what actually holds — and
declared bounds are harmless for the inclusion, because every runtime
instance was constructed through a \texttt{new}-check that enforced them
(under \textbf{L-TC} at construction sites).

\subsection{Conditional discharge and the honest caveats}

\begin{theorem}[Conditional discharge of \textbf{L-INST}]
\label{thm:inst}
Assume, for the class $C$ at a narrowing site:
\begin{itemize}
\item \textbf{A-HASINST}: no \texttt{@@hasInstance} method other than
  the \texttt{Function.prototype} default is found on
  $\realize(C)$ (step 2 resolves to the default);
\item \textbf{A-PROTO}: no prototype surgery — objects with
  $C.\texttt{prototype}$ in their chain were produced by \texttt{new} on
  $C$ or a subclass (no \texttt{Object.create},
  \texttt{Object.setPrototypeOf}, \texttt{Reflect}, or Proxy
  $\intm{GetPrototypeOf}$ traps reachable), and
  $C.\texttt{prototype}$ itself is unmodified;
\item \textbf{L-TC} at construction: every \texttt{new}-site
  establishes the instance object type of the constructed class
  (fields present and typed, per the class-checking discipline).
\end{itemize}
Then $P_{\mathsf{inst}}(C) \subseteq \sem{T_C}_{\rho'}$ at the
post-instantiation state $\rho'$, reading the fresh unbound arguments'
denotations as unconstrained-within-declared-bounds.
\end{theorem}
\begin{proof}[Proof sketch]
By \textbf{A-HASINST}, membership in $P_{\mathsf{inst}}(C)$ is decided by
\texttt{OrdinaryHasInstance}, so $v$ is an object with
$C.\texttt{prototype}$ in its chain. By \textbf{A-PROTO}, $v$ was
constructed by $C$ or a subclass, as some instantiation
$C\langle \vec U \rangle$ with $\vec U$ within the declared bounds. By
\textbf{L-TC} at the construction site, $v$ then satisfies the instance
object type of $C\langle \vec U \rangle$: every declared field present
with a $\vec U$-instantiated get-type, inherited fields included
(subclass constructors must complete \texttt{super()}). Since $T_C$'s
arguments are \emph{fresh and unbound at $\rho'$}, each field constraint
of $T_C$ weakens to its declared bound, which $\vec U$ satisfies; so
$v \in \sem{T_C}_{\rho'}$. Mutation after construction is governed by
\textbf{L-MUT} exactly as for every other membership statement in the
main paper.
\end{proof}

Each hypothesis must now be audited against the implementation. The
audit \emph{fails} for the first, and the aftermath of the theorem
fails for the freshness clause; both failures are pinned.

\begin{finding}[\textbf{A-HASINST} is violated: SKJS accepts
\texttt{@@hasInstance} overrides and narrows anyway]
\label{find:hasinstance}
The narrowing gate checks only that the RHS name has a class definition
(\texttt{hasClassDef}); nothing rejects or inspects a static
\texttt{[Symbol.hasInstance]} member, which SKJS class bodies \emph{can}
declare — computed symbol-keyed members resolve into symbol fields
(the \texttt{FNComputed} arm of \texttt{addFields}), and
\texttt{Symbol.hasInstance} is a stubbed \texttt{unique symbol}.
The implementation itself documents the gap
(``Mirrors tsc's TS2359 acceptance (modulo \texttt{Symbol.hasInstance},
which SKJS does not model)'', \texttt{isInstanceofRhs}). Checked
witness (accepted at the pin, exit 0):
\begin{verbatim}
class C {
  a: number = 1;
  static [Symbol.hasInstance](x: unknown): boolean { return true; }
}
function f(x: unknown): number {
  if (x instanceof C) { return x.a; }   // narrowed to C
  return 0;
}
console.log(f("not a C"));              // Node prints: undefined
\end{verbatim}
At runtime step 2 of \texttt{InstanceofOperator} finds the static
method, every value is ``an instance'', and \texttt{f} returns
\texttt{undefined} at type \texttt{number}. Two mitigations narrow the
blast radius: (i) the override must be declared \emph{in the checked
program itself} on a class the same program narrows on — ambient code
cannot smuggle one in unseen, since the stubs declare none; (ii) the
hole is shared with \texttt{tsc}, which also narrows here (a
boolean-returning \texttt{@@hasInstance} does not disable
constructor-based narrowing there; only \texttt{is}-predicate returns
change it). Cheap repairs exist: reject static
\texttt{[Symbol.hasInstance]} declarations, or skip narrowing for
classes that declare one (the definition is available at the gate).
\impl{TypeCheckRefine.sk / evalCond (InstanceOf arm)}
\impl{TEnv.sk / addFields (FNComputed), resolveComputedKeyFieldKey}
\impl{TypeCheckExpr.sk / isInstanceofRhs (comment)}
\end{finding}

\begin{finding}[Instantiation weakness is real: fresh arguments are
sealed by later covariant use]
\label{find:fresh-seal}
Theorem~\ref{thm:inst} holds \emph{at the narrowing instant}, and the
main paper's \textsc{R-Inst} lemma correctly delegates the survival of
membership under later substitution evolution to its
Obligation~\emph{Monotone Reads Under Growth} — whose status there is
\textsc{violated}, this finding being one of its two witnessed
instances:
a later covariant read of an unbound candidate does not seal it to
$\top$ but to the \emph{expected type of the use}
(the $\mathsf{grow}(t, \mathsf{cov})$/unbound $\to$
$\mathsf{SVClosed}(t)$ transition of the main paper's Definition
``Transitions''), so
program \emph{uses} choose the type argument that narrowing had left
open — the unsound direction, since no value flow justifies it. Checked
witness (accepted at the pin, exit 0). \texttt{tsc} 5.7 also accepts
this program, but through its declared \texttt{any} escape hatch: it
narrows \texttt{unknown} to \texttt{Box<any>} and accepts a
\texttt{string} and a \texttt{number} read simultaneously. SKJS's
failure mode is different and worse — \texttt{any}-free and
order-dependent: the first covariant use wins, and a second,
differently-typed read is rejected pointing at the first. The witness:
\begin{verbatim}
class Box<T> { v: T; constructor(v: T) { this.v = v } }
function h(x: unknown): string {
  if (x instanceof Box) {
    const s: string = x.v;   // seals the fresh argument to string
    return s;
  }
  return "no";
}
console.log(h(new Box(42)));  // Node prints: 42
\end{verbatim}
$h$ returns \texttt{42} at type \texttt{string}. The \texttt{Array}
special case ($\top$ argument) is immune, which suggests the repair:
instantiate \emph{all} \texttt{instanceof} type arguments at $\top$
(an explicit top argument — matching tsc in \emph{effect} while avoiding
its \texttt{any} escape hatch), or mark
narrowing-born variables as non-sealable-by-use. Until then, the
freshness clause of \textbf{L-INST}'s discharge note is \emph{not
sufficient}, and the violated instance of the monotonicity obligation
should be recorded in the ledger (Section~\ref{sec:status}).
\impl{TEnv.sk / makeTArgs; TTypeSubst.sk / Subst::grow (cov/unbound)}
\impl{TypeCheckRefine.sk / evalInstanceof (Array $\top$ special case)}
\end{finding}

\paragraph{\textbf{A-PROTO} audit.}
The pinned stubs omit the surgery surface deliberately: the
\texttt{Object} statics \texttt{create}, \texttt{getPrototypeOf}, and
\texttt{setPrototypeOf} are present only as comments marked
``disallowed'', and neither \texttt{Proxy} nor \texttt{Reflect} is
declared at all. A checked program therefore cannot express prototype
surgery; what remains is the generic unchecked boundary (ambient
\texttt{declare}d functions, \texttt{JSON.parse<T>}) through which any
assumption, not specially this one, can be voided — that boundary is
priced into \textbf{L-TC} already. Status: \textsc{assumed}, adequately
enforced for checked code.
\impl{import-stubs/@types/skjs/index.d.ts (Object statics)}

\paragraph{Cross-realm objects.}
\texttt{instanceof} is realm-sensitive: an array (or class instance)
from another realm does not have \emph{this} realm's
\texttt{C.prototype} in its chain, so the operator answers
\textsf{false}. For soundness of \textsc{R-Inst} this is \emph{benign}:
the rule narrows only the true branch, and cross-realm values make the
test false, landing in the untouched (identity) false branch. Realms are
purely a \emph{completeness} phenomenon for \texttt{instanceof} — the
classical reason \texttt{Array.isArray} exists, as the next section
exploits. No realm assumption is needed for \textbf{L-INST}; we record
this so the ledger does not grow a vacuous row.

%==========================================================================
\section{\texttt{Array.isArray} and \textbf{L-ISARR}}
\label{sec:isarr}

\subsection{The spec side}

\texttt{Array.isArray} (\ecma{23.1.2.2}) is \texttt{IsArray}
(\ecma{7.2.2}): return \textsf{true} if the argument is an
\emph{Array exotic object} (\ecma{10.4.2}: the objects with the special
\texttt{length}-coupled $\intm{DefineOwnProperty}$); if it is a Proxy,
recurse on the Proxy's target (throwing a \texttt{TypeError} if the
proxy has been revoked); otherwise \textsf{false}. Three consequences:
(i) the test is \emph{cross-realm-safe} — an array from any realm is an
Array exotic object, no prototype identity involved — in contrast to
\texttt{x instanceof Array}; (ii) it is \emph{Proxy-transparent} — a
proxy over an array answers \textsf{true}; (iii) on a revoked proxy the
\emph{call throws}, so by the main paper's evaluation convention neither
branch relation holds and narrowing is vacuously sound for that case.
Ground $\mathsf{isArrayV}(o) \triangleq \texttt{IsArray}(\realize(o))$
(defined, non-throwing case).

\subsection{The model side: how \texttt{ReadonlyArray} acquires the
brand}

\textbf{L-ISARR} claims $\mathsf{isArrayV}(o)$ implies
$o \in \sem{\mathsf{ReadonlyArray}\langle\top\rangle}$. Two things could
break this: the stub type could demand a \emph{property} real arrays do
not have, or a declared member could fail on real arrays. The first
danger is real on the surface: the stub spells the brand as an ordinary
declaration,
\begin{verbatim}
declare interface ReadonlyArray<T> extends Iterable<T> {
  readonly __SKJS_isArray: true;
  ...
}
\end{verbatim}
and a required field \texttt{\_\_SKJS\_isArray} would exclude every real
runtime array (they carry no such property), making
$\sem{\mathsf{ReadonlyArray}}$ nearly empty and \textbf{L-ISARR} false.
The implementation avoids this by construction: \texttt{TFields::create}
\emph{intercepts} the field named by
\texttt{kArrayBrandFieldName = "\_\_SKJS\_isArray"}, removes it from the
field map, and transfers its literal value (\texttt{true}/\texttt{false},
anything else is a fatal SKJS2051) into the dedicated
\texttt{TFields.isArray :\ ?Bool} component — the model's brand
$\mathsf{isArr}$, whose membership semantics constrains
$\mathsf{isArrayV}$, not $\mathsf{props}$. The brand never exists as an
ordinary field of any type built through \texttt{TFields::create}, which
is the path all class/interface instantiation takes
(\texttt{TEnv::newObject}). Defaults are set soundly at the same spot:
interfaces get $\mathsf{isArr} = \mathsf{None}$ (unknown) and classes
get $\mathsf{Some}(\mathsf{false})$ \emph{before} the intercept — so a
class inheriting the brand field (e.g.\ \texttt{Array} itself, tuples
via \texttt{SKJS\_Tuple}, or a user \texttt{class D extends Array})
still ends at $\mathsf{Some}(\mathsf{true})$, because the inherited
brand field is in the collected map and the intercept overrides the
default. (Object \emph{literals} take \texttt{createFromSorted}, which
does not intercept; a user field literally named
\texttt{\_\_SKJS\_isArray} on a literal stays an ordinary — truthfully
present — property, which is sound.)
\impl{TType.sk / kArrayBrandFieldName, TFields::create (brand
intercept)}
\impl{TEnv.sk / newObject (isArray defaults)}
\impl{import-stubs/@types/skjs/index.d.ts / ReadonlyArray}

\subsection{Discharge}

\begin{theorem}[Discharge of \textbf{L-ISARR}]
\label{thm:isarr}
Assume \textbf{A-PROXY}: no Proxy exotic object flows into checked code.
Then $\mathsf{isArrayV}(o)$ implies
$o \in \sem{\mathsf{ReadonlyArray}\langle\top\rangle}$.
\end{theorem}
\begin{proof}
Under \textbf{A-PROXY}, $\realize(o)$ is an Array exotic object of some
realm. $\sem{\mathsf{ReadonlyArray}\langle\top\rangle}$ has a single
object branch whose constraints we check against the object interface
grounded as: $\mathsf{has} =$ \texttt{HasProperty} (prototype chain
included — this matches both the \texttt{in} operator and property
reads) and $\mathsf{get} = \intm{Get}$, snapshots per \textbf{L-PURE}.
\begin{itemize}
\item \emph{Brand}: $\mathsf{isArr} = \mathsf{Some}(\mathsf{true})$
  (intercepted as above) requires $\mathsf{isArrayV}(o)$ — the
  hypothesis.
\item \emph{\texttt{length}}: required, get-type \texttt{number}. Every
  Array exotic object owns a \texttt{length} whose value is a
  non-negative Number (\ecma{10.4.2}).
\item \emph{Numeric index} (\texttt{readonly [n: number]:\ T} at
  $T = \top$): covers numeric-spelling keys not otherwise named; each
  such property value lies in $\sem{\top} = \Val$. Trivially satisfied
  (and the $\mathsf{set} = \bot$ readonly encoding is invisible to
  membership, which is get-only).
\item \emph{Methods} (\texttt{find}, \texttt{every}, \texttt{some},
  \texttt{forEach}, \texttt{findIndex}, \texttt{join}, \texttt{map},
  \texttt{filter}, \texttt{flatMap}, \texttt{reduce},
  \texttt{reduceRight}, \texttt{slice} — the full declared list at the
  pin, several classic members being deliberately commented out for
  variance reasons): each is a required field whose get-type is a
  single-branch object with a \emph{definite} call signature
  (interface methods are built through \texttt{TObject::createCall}).
  Membership requires presence and callability of the property value;
  the behavioral meaning of the signature itself is opaque
  (bundled into the subtyping surface by the main paper's object semantics).
  Every listed name is a function-valued property of
  \texttt{Array.prototype} in every realm (\ecma{23.1.3}), reached by
  $\mathsf{get}$ through the instance's prototype chain; arrays from
  other realms find their own realm's identical-by-spec
  \texttt{Array.prototype}. (An own property shadowing a method with a
  non-callable value would falsify this — that is ordinary property
  mutation, priced into \textbf{L-MUT}/\textbf{L-PURE} snapshots like
  every other field statement, not a new caveat.)
\item \emph{\texttt{[Symbol.iterator]}} (inherited from the
  \texttt{Iterable<T>} base): same argument; the \texttt{@@iterator}
  method of \texttt{Array.prototype} exists in every realm, and
  well-known symbols are \emph{shared across realms}
  (\ecma{6.1.5.1}), so the checker's unique-symbol key denotes the same
  runtime symbol everywhere.
\item \emph{Extensibility}: all other properties of $o$ are
  unconstrained by the main paper's extensible object semantics; no
  string index is declared, so no covering constraint arises.
\end{itemize}
All constraints hold; $o$ inhabits the branch.
\end{proof}

On \textbf{A-PROXY}: a Proxy over an array passes \texttt{IsArray} but
its $\intm{Get}$ trap may return arbitrary values for \texttt{length} or
the methods, falsifying membership; hence the assumption. It is enforced
for checked code the same way as \textbf{A-PROTO}: \texttt{Proxy} is
absent from the pinned stubs, so no checked program can construct one,
and the unchecked boundary is \textbf{L-TC}'s. Note the contrast worth
keeping in the ledger note: cross-realm values, the classical failure
mode of \texttt{instanceof Array}, are handled \emph{correctly} here —
no realm assumption is required, thanks to \texttt{IsArray}'s
realm-independence plus the structural (non-nominal) membership
semantics.

The mutability note of the main paper (\textsc{R-IsArray} intersects
with \emph{Readonly}Array, never mutable \texttt{Array}) needs no
runtime grounding: it is a precision choice about which \emph{set-types}
to add, and set-types do not participate in membership.

%==========================================================================
\section{Status and Ledger Edits}
\label{sec:status}

``Here'' means in this deep dive.

\begingroup
\renewcommand{\arraystretch}{1.25}
\small
\begin{longtable}{@{}>{\raggedright\arraybackslash}p{4.2cm}>{\raggedright\arraybackslash}p{9.2cm}@{}}
\toprule
\textbf{Item} & \textbf{Where} \\
\midrule
\endfirsthead
\toprule
\textbf{Item} & \textbf{Where} \\
\midrule
\endhead
\midrule
\multicolumn{2}{r@{}}{\small\emph{Continued on next page}} \\
\endfoot
\bottomrule
\endlastfoot
L-TAG & Theorem~\ref{thm:tags}\\
$\approx$ grounding & Theorem~\ref{thm:seq},
  Corollary~\ref{cor:zero}\\
Truthiness & Theorem~\ref{thm:truthy}\\
ctor$\Rightarrow$call & Theorem~\ref{thm:ctor-call}\\
\textbf{A-DDA} & \S\ref{sec:typeof-table}\\
\textbf{A-CTOR-CALL} & Finding~\ref{find:ctor-call}\\
L-INST & Theorem~\ref{thm:inst}\\
\textbf{A-HASINST} & Finding~\ref{find:hasinstance}\\
\textbf{A-PROTO} & \S\ref{sec:inst}\\
Fresh-arg survival & Finding~\ref{find:fresh-seal}\\
L-ISARR & Theorem~\ref{thm:isarr}\\
\textbf{A-PROXY} & \S\ref{sec:isarr}\\
\end{longtable}
\endgroup

\paragraph{Registry.}
The authoritative, current statuses of every item above are maintained in
the registry (\texttt{docs/ledger.yaml}), rendered as the main paper's
Assumptions and Proof Obligations tables; this paper is cited there as
evidence. The registry rows it supports:

% GENERATED by docs/tools/docs.py from docs/ledger.yaml — do not edit.

\begin{center}
\renewcommand{\arraystretch}{1.25}
\small
\begin{tabular}{@{}>{\raggedright\arraybackslash}p{2.4cm}>{\raggedright\arraybackslash}p{8.2cm}>{\raggedright\arraybackslash}p{2.6cm}@{}}
\toprule
\textbf{Row} & \textbf{Statement} & \textbf{Status} \\
\midrule
L-TAG & Every runtime value carries exactly one tag consistent with its static tag. & \textsc{proved} \\
A-CTOR-CALL & Class-constructor values are typed with call mode $\mathsf{None}$, but constructors implement $[[\mathrm{Call}]]$. & \textsc{violated} \\
L-INST & $P_{\mathsf{inst}}(C) \subseteq \sem{T_C}$ for the instance type as instantiated at the narrowing site. & \textsc{violated} \\
L-ISARR & $\mathsf{isArrayV}(o)$ implies $o \in \sem{\mathsf{ReadonlyArray}\langle\top\rangle}$. & \textsc{proved} \\
A-PROTO & No prototype surgery / Proxy GetPrototypeOf traps reach checked code (enforced by stub omissions). & \textsc{assumed} \\
A-DDA & No [[IsHTMLDDA]] object (\texttt{document.all}) reaches checked code (vacuous on Node). & \textsc{assumed} \\
A-PROXY & No Proxy exotic object reaches checked code. & \textsc{assumed} \\
A-HASINST & No non-default \texttt{Symbol.hasInstance} on narrowed \texttt{instanceof} RHS classes. & \textsc{violated} \\
\bottomrule
\end{tabular}
\end{center}


\paragraph{Imports (review contract).}
The named statements this paper assumes from elsewhere. A reviewer treats
each as an axiom, checking citations against the registry's canonical
statement; no other paper's internals are in scope:

% GENERATED by docs/tools/docs.py from docs/ledger.yaml — do not edit.

\begin{center}
\renewcommand{\arraystretch}{1.25}
\small
\begin{tabular}{@{}>{\raggedright\arraybackslash}p{2.6cm}>{\raggedright\arraybackslash}p{7.4cm}>{\raggedright\arraybackslash}p{1.2cm}>{\raggedright\arraybackslash}p{2.2cm}@{}}
\toprule
\textbf{Row} & \textbf{Canonical statement} & \textbf{Owner} & \textbf{Status} \\
\midrule
L-INTER & When \texttt{inter} returns a non-$\bot$ result, that result contains the semantic intersection of its operands, provided every recursive call that returned $\bot$ had disjoint operands (L-INTER-BOT). & B & \textsc{open} \\
L-INTER-BOT & A $\bot$ result from \texttt{inter} implies semantic disjointness. & B & \textsc{violated} \\
L-TC & Checked expressions inhabit their checked types. & main & \textsc{assumed} \\
L-PURE & Evaluating a condition does not mutate the narrowed locals. & F & \textsc{violated} \\
L-MUT & Field writes are governed by $\mathsf{set}$-types; get-only membership is sound for snapshot reads. & main & \textsc{assumed} \\
L-MUT-ALIAS & A checked field or index write through one mutable view preserves the membership guarantees of every live alias of the same runtime object. & H & \textsc{violated} \\
OBL-MONOTONE & Monotone Reads Under Growth: judgements validated under $\rho$ survive later legal transitions. & E & \textsc{violated} \\
TVAR-SEM & A type-variable occurrence denotes its substitution envelope intersected with its occurrence constraints; an unbound envelope is empty. & main & \textsc{proved} \\
SKOLEM-SEM & A skolem in an installed \texttt{instanceof} view denotes the scrutinee object's actual hidden type argument, fixed at narrowing time and stable thereafter. & main & \textsc{proved} \\
F-CABSENT & The $\mathsf{CAbsent}$ brand-like claim accurately excludes callables: refinement drops $\mathsf{CAbsent}$ branches under \texttt{typeof} \texttt{"function"}, but subtyping's permissive $(\mathsf{CUnknown},\mathsf{CAbsent})$ arm admits callables into branded types. & D & \textsc{violated} \\
\bottomrule
\end{tabular}
\end{center}


\paragraph{Method note (failed attempts).}
Three discharge attempts failed and became the Findings above; each
failure was converted into a checked-and-executed witness rather than a
weakened assumption, per the series contract. One further attempt is
worth recording as simply \emph{unnecessary}: we initially looked for a
realm assumption for \textbf{L-INST} symmetric to the classical
\texttt{instanceof Array} pitfall, and found that true-branch-only
narrowing makes realms a pure completeness concern
(\S\ref{sec:inst}), so no such row exists.

%==========================================================================
\appendix
\section{Implementation and Specification Anchors}
\label{app:anchors}

All implementation references are to revision \texttt{9108f339}.

\begin{center}
\renewcommand{\arraystretch}{1.25}
\small
\begin{tabular}{@{}l@{\quad}l@{}}
\toprule
\textbf{This paper} & \textbf{Anchor} \\
\midrule
brand name & \texttt{TType.sk / kArrayBrandFieldName} \\
brand intercept & \texttt{TType.sk / TFields::create} \\
brand defaults & \texttt{TEnv.sk / newObject} \\
brand stub spelling & \texttt{import-stubs/@types/skjs/index.d.ts} \\
\textsc{R-IsArray} handler & \texttt{TypeCheckRefine.sk / evalIsArray} \\
array-brand intersection & \texttt{TTypeInter.sk /
  TFields::interIsArray} \\
\textsc{R-Inst} gate + handler & \texttt{TypeCheckRefine.sk / evalCond,
  evalInstanceof} \\
instance type arguments & \texttt{TEnv.sk / makeTArgs};
  \texttt{TTypeSubst.sk / Subst::gen} \\
class-def gate & \texttt{TEnv.sk / hasClassDef} \\
RHS validity (TS2359) & \texttt{TypeCheckExpr.sk / isInstanceofRhs} \\
symbol-keyed class members & \texttt{TEnv.sk / addFields (FNComputed)} \\
static-side class type & \texttt{TEnv.sk / newStaticObject} \\
\textsc{R-Typeof} handler & \texttt{TypeCheckRefine.sk /
  evalTypeofEquality} \\
truthiness tables & \texttt{TType.sk / TPrimValueKind::isFalsy,
  SType::falsiness} \\
\bottomrule
\end{tabular}
\end{center}

Specification clauses cited (ECMA-262, 2025-edition numbering; clause
names are stable across editions): \S6.1 language types; \S6.1.5.1
well-known symbols; \S6.1.6.1--6.1.6.2 Number/BigInt types (and
\texttt{Number::equal}, \S6.1.6.1.13); \S6.1.7.2--6.1.7.3 object
internal methods and their invariants; \S7.1.2 \texttt{ToBoolean};
\S7.2.2 \texttt{IsArray}; \S7.2.3--7.2.4 \texttt{IsCallable} and
\texttt{IsConstructor}; \S7.2.10--7.2.12 \texttt{SameValue},
\texttt{SameValueZero}, and \texttt{SameValueNonNumber};
\S7.2.16 \texttt{IsStrictlyEqual}; \S7.3.22
\texttt{OrdinaryHasInstance}; \S10.2.1 ECMAScript function
$\intm{Call}$ (class-constructor throw); \S10.4.1 bound function exotic
objects; \S10.4.2 Array exotic objects; \S10.5.12--10.5.13 Proxy
$\intm{Call}$ and $\intm{Construct}$; \S13.5.3 \texttt{typeof};
\S13.10.2 \texttt{InstanceofOperator}; \S13.11 equality operators;
\S20.2.3.6 the \texttt{@@hasInstance} method of
\texttt{Function.prototype}; \S23.1.2.2
\texttt{Array.isArray}; \S23.1.3 \texttt{Array.prototype} methods;
Annex~B \S B.3.6 $\intm{IsHTMLDDA}$.

\section{Witness Protocol}
\label{app:witness}

Each Finding's program was (i) written to a temporary file in the
repository working tree, (ii) checked with the pinned binary — the
\texttt{make check-file-in-docker} target, passing
\texttt{BIN=target/host/dev/skjs} — all three accepted with
exit code 0
and no diagnostics — and (iii) type-stripped mechanically and executed
with Node to observe the printed value named in the Finding
(\texttt{[class C]}, \texttt{undefined}, \texttt{42} respectively). The
temporary files were removed afterwards; turning the witnesses into
pinned tests is proposed in Section~\ref{sec:status} but deliberately
not done here (this deep dive changes no file outside
\texttt{docs/deep-dives/}).

\end{document}
